\documentclass[11pt]{article}

\usepackage[margin=1in]{geometry}
\usepackage{amsmath}
\usepackage{amssymb}
\usepackage{booktabs}
\usepackage{graphicx}
\usepackage{hyperref}
\usepackage{authblk}
\usepackage{xcolor}
\usepackage{caption}
\usepackage{microtype}
\usepackage[round]{natbib}

\hypersetup{
    colorlinks=true,
    linkcolor=black,
    citecolor=black,
    urlcolor=blue!60!black,
}

\title{Statistical Arbitrage in Crypto Pair Spreads:\\
Kalman Hedging, Selection Bias, and the Limits of Predictability\\
\large{Phase 1 Report}}

\author{Jack Lutz \and Sammy Adham \and Frank Kronewitter}
\affil{University of California, Los Angeles}
\date{\today}

\begin{document}
\maketitle

\begin{abstract}
We study statistical arbitrage on hourly Binance.US USDT spot pairs,
adapting the equity pairs-trading framework of \citet{gatev2006pairs} to
cryptocurrency markets. The central methodological finding is that a static
ordinary-least-squares (OLS) hedge ratio fails out-of-sample on this
market: of 1{,}219 candidate pairs in the top-50 liquidity universe at
$t_0 = 2024\text{-}01\text{-}01$, only $5.4\%$ maintain an augmented
Dickey--Fuller (ADF) $p < 0.05$ on a held-out 30-day test slice when the
hedge ratio is fit only on the preceding 90-day training window. A Kalman
filter with parameters maximum-likelihood-fit on training data and
forward-rolled on test recovers cointegration on $99.7\%$ of the same
pairs at the same threshold. The relationships in crypto pairs are real
but non-stationary; a fixed hedge ratio is the wrong model.

We further evaluate four classifiers (a gradient-boosted tree, an LSTM,
and two transformers of different size) on a 27-split walk-forward
protocol across the Kalman-derived spreads. The classifiers' predictive
signal collapses, in a one-feature-out ablation, almost entirely to the
single feature $\text{spread}_z$, which is the same quantity the
classical $z$-threshold rule already uses. After a Bailey--L\'opez de
Prado deflated-Sharpe correction for the $N \approx 50$ strategy
configurations evaluated in this project, the pre-cost annualized Sharpe
of $\sim$8 fails to reject the null of zero true Sharpe at realistic
trial counts; at maker-fee-equivalent transaction costs the realized
Sharpe drops to $\sim$0.9 and sits below the chance-best of $\sim$1.2 at
$N = 10$. We retract the corresponding break-even claim.

The narrower defensible market-efficiency statement: at hourly cadence
on Binance.US USDT pairs, the mean-reverting alpha resides in the
$|\text{spread}_z| \geq 2$ entry condition itself, not in the choice of
classifier. The Kalman-spread cointegration result is robust and the
load-bearing finding of this phase. L2 microstructure data, expected in
Phase~2, is required to test whether intra-bar execution improves the
post-cost picture.
\end{abstract}

\section{Introduction}

\citet{gatev2006pairs} (henceforth GGR) propose a systematic
pairs-trading rule on US equities: select pairs by historical distance,
trade the spread when it exceeds $2$ standard deviations from its mean,
close at convergence. Their reported Sharpe of $\sim$0.5--1.0 has
motivated decades of follow-on work in equities, commodities, and more
recently digital assets.

The cryptocurrency market differs from the GGR setting in three
material ways. First, the investable universe changes rapidly: new
tokens are listed weekly and others are delisted, often without warning.
Second, the typical relationship between two tokens is unstable on
multi-month horizons; pairs that look cointegrated on one year of data
may show no relationship the next. Third, exchange fees are large
relative to typical spread movements: Binance.US default taker fees of
10 bps per side imply 20 bps per round-trip per leg of a pair trade.

This work adapts the GGR framework to hourly Binance.US spot data with
USDT quote, with two specific contributions:

\begin{enumerate}
    \item A direct out-of-sample comparison of static OLS hedge ratios
    against a Kalman-filter dynamic hedge across the full universe of
    1{,}219 liquidity-filtered pairs, showing that the time-varying
    hedge recovers cointegration almost universally where the static
    hedge fails.
    \item An honest evaluation, using block-bootstrap confidence
    intervals and a deflated Sharpe correction \citep{bailey2014dsr},
    of whether modern machine-learning models provide a tradable edge
    on Kalman-derived spreads after transaction costs and after
    accounting for the multiple-testing inherent in our methodology
    sweep.
\end{enumerate}

We find that the first contribution is unambiguous and constitutes the
load-bearing finding of this phase. The second is largely negative: the
ML models we tried do not produce a selection-corrected Sharpe that
beats chance at realistic transaction-cost levels, and most of the
``alpha'' visible in our backtests is attributable to the
$|z| \geq 2$ entry condition rather than the classifier that fires it.

The remainder of this paper is organized as follows. Section~\ref{sec:data}
describes the data and the as-of-time universe construction. Section~\ref{sec:methods}
details the methodology, including the Kalman state-space specification,
the feature set, the classifier architectures, and the evaluation
protocol. Section~\ref{sec:results} presents the four headline findings.
Section~\ref{sec:limits} catalogs limitations. Section~\ref{sec:phase2}
describes the Phase~2 L2-microstructure plan.

\section{Data} \label{sec:data}

\paragraph{Source.} Hourly open-high-low-close-volume (OHLCV) candles
from the Binance.US REST endpoint \texttt{api.binance.us/api/v3/klines},
USDT quote only. The local store contains 195 currently-listed symbols,
with first-observation dates ranging from 2019-09-23 (BTC, ETH, BNB,
XRP, BCH) to 2025-12-18. All symbols share the same final bar at
2026-01-22 04:00 UTC, the snapshot date.

\paragraph{As-of-time universe.} A symbol enters the universe at
calendar time $t_0$ if and only if its first observation precedes
$t_0 - \tau_{\min}$ and its last observation postdates $t_0 - \Delta t$,
where $\tau_{\min}$ is a minimum-history requirement and $\Delta t$ is
one bar. We use $\tau_{\min} = 180$ days throughout. The first-observation
check is necessary to prevent look-ahead leakage when $t_0$ is in the
past; an earlier version of our universe function omitted this check
and included symbols whose listing postdated the analysis cutoff.

\paragraph{Liquidity filter.} Within the as-of universe at $t_0$, we
rank symbols by mean USDT volume ($\text{close} \times \text{volume}$)
over the training window and keep the top 50. Symbols with less than
$80\%$ coverage of the expected hourly calendar in the window are
dropped before ranking. At $t_0 = 2024\text{-}01\text{-}01$, this
yields a 130-symbol universe and a top-50 liquid set.

\paragraph{Survivorship.} Because the snapshot was downloaded as a
single point in time and only contains currently-listed symbols, the
panel cannot observe pre-snapshot delistings. We compiled an external
delisting log from the Binance.US Help Center announcements: 14 tokens
were delisted between $t_0$ and the snapshot date.\footnote{BOND, ANT,
WAVES, TUSD, CUDOS, MXC, REN, VITE, BAL, CLV, STMX, LOOM, KDA, OXT, JAM.}
This implies a survivorship rate of approximately $14 / 209 = 6.7\%$ as
a loose upper bound. After the liquidity top-50 filter, only an
estimated 3--6 of these would have been candidates; the remainder were
small-cap tokens that would not have passed the liquidity threshold
even at $t_0$. The realistic bias on the headline metrics is therefore
$5\%$--$12\%$ of the pre-filter candidate set.

\section{Methodology} \label{sec:methods}

\subsection{Pair selection}

We screen pairs in two stages. The strict cointegration screen, adapted
from GGR, requires (i) Pearson correlation of log returns $\geq 0.5$,
(ii) ADF $p$-value on the OLS residual $\leq 0.05$, (iii) Ornstein--Uhlenbeck
half-life in $[4, 240]$ hours, and (iv) beta stability (standard deviation
of OLS $\beta$ across four equal segments, divided by absolute mean
$\beta$) $\leq 0.35$. On the local Binance.US slice at
$t_0 = 2024\text{-}01\text{-}01$, this strict screen finds \emph{zero}
passing pairs. We therefore use a correlation fallback (rank by Pearson
correlation of returns, take top $N$) for the downstream pipeline, and
treat the failure of the strict screen as evidence that motivates the
Kalman analysis of Section~\ref{sec:kalman}.

\subsection{Static and dynamic hedge ratios} \label{sec:kalman}

For each pair $(A, B)$ with log prices $y_t = \log p^A_t$ and
$x_t = \log p^B_t$, the \emph{static OLS hedge} is the pair
$(\alpha, \beta)$ minimizing $\sum_t (y_t - \alpha - \beta x_t)^2$ over
a training window. The residual $\epsilon_t = y_t - \alpha - \beta x_t$
is the static spread.

The \emph{Kalman dynamic hedge} treats $(\alpha, \beta)$ as a hidden
state evolving as a random walk:
\begin{align}
    \mathbf{s}_t &= \begin{bmatrix} \alpha_t \\ \beta_t \end{bmatrix},
    \qquad
    \mathbf{s}_t = \mathbf{s}_{t-1} + \mathbf{w}_t,
    \qquad \mathbf{w}_t \sim \mathcal{N}(\mathbf{0}, Q), \label{eq:state}\\
    y_t &= \begin{bmatrix} 1 & x_t \end{bmatrix} \mathbf{s}_t + v_t,
    \qquad v_t \sim \mathcal{N}(0, R), \label{eq:obs}
\end{align}
with $Q = \mathrm{diag}(q_\alpha, q_\beta)$. The process noise $Q$ and
observation noise $R$ are unknown parameters and are fit by maximum
likelihood on a training slice, with the likelihood being the standard
prediction-error decomposition,
\begin{equation}
    \log L(Q, R) = -\tfrac{1}{2} \sum_t \left[ \log(2\pi S_t)
    + \frac{(y_t - \hat{y}_{t|t-1})^2}{S_t} \right],
\end{equation}
where $\hat{y}_{t|t-1}$ is the one-step-ahead prediction and $S_t$ is
its variance under the filter. We optimize $(\log q_\alpha, \log q_\beta,
\log R)$ by L-BFGS-B; the log parameterization keeps the variances
positive.

After fitting, we run the filter forward once more on the training slice
with the fitted parameters, record the final filtered state
$\mathbf{s}_{T_{\text{train}}}$ and the training-period residuals
$\hat{\epsilon}_t^{\text{train}}$. On the test slice we initialize the
filter at $\mathbf{s}_{T_{\text{train}}}$, hold the fitted $(Q, R)$
fixed, and produce out-of-sample residuals
$\hat{\epsilon}_t^{\text{test}}$. We perform \emph{no} re-fitting on
test data. The training and test residual series are concatenated into
the spread series used downstream, with the same filter parameters
operative on both sides of the train/test boundary.\footnote{An earlier
version of this pipeline used different parameters on the two sides
(default Kalman config for training, MLE-fit for test), which produced
a discontinuity at the boundary. This was caught in audit and fixed;
see Section~\ref{sec:audit}.}

\subsection{Features and labels}

Per pair, per hourly bar, we compute 16 features in three groups:
spread features ($\text{spread}$, $\text{spread}_z^{168h}$,
$\text{spread}_{\Delta}^{1h, 4h, 24h}$, $\text{spread}_{\sigma}^{24h, 168h}$),
volume features ($\text{volume ratio}$, $\text{volume }z\text{-scores}$,
$\text{pair volume sum/geometric mean}$), and market-context features
($\text{BTC return}^{24h}$, $\text{rolling correlation}^{24h}$,
$\text{realized volatility}^{24h}_{A, B}$). All features are
backward-rolling. The target is the change in absolute spread over the
next 24 hours,
\begin{equation}
    \Delta_t = |s_{t+24}| - |s_t|,
\end{equation}
quantized into three classes: $0$ (revert, $\Delta_t < -\tau$), $1$
(persist, $|\Delta_t| \leq \tau$), $2$ (diverge, $\Delta_t > \tau$).
The threshold $\tau$ is either a fixed scalar (0.001 in log-spread
units) or a per-pair scalar
$\tau_{\text{pair}} = c \cdot \widehat{\sigma}(\Delta_t)_{\text{pair}}$
computed from the pair's training slice only. The per-pair option
normalizes label distributions across pairs of different spread
volatility. Throughout the experiments below, $\tau$ is set per-pair
with $c = 0.5$ and the per-pair training slice ends at $t_0$.

\subsection{Models}

We compare seven baselines and learned classifiers on the same
sliding-window dataset (window length 168 hours, stride 1 hour,
horizon 24 hours):

\begin{itemize}
    \item \texttt{persist\_class}: always predict 1 (no trade).
    \item \texttt{majority\_class}: predict the training mode.
    \item \texttt{random\_stratified}: sample classes with training prior.
    \item \texttt{zscore\_rule}: predict 0 (revert) if
    $|\text{spread}_z| \geq 1.5$, else 1.
    \item \texttt{booster}: \texttt{sklearn HistGradientBoosting},
    300 trees, learning rate $0.05$, on the tabular feature summary
    of each 168-hour window.
    \item \texttt{lstm}: 2-layer LSTM, hidden dimension 64, dropout 0.2,
    linear head to 3 logits, trained 4 epochs with Adam and learning
    rate $10^{-3}$.
    \item \texttt{transformer (small)}: 2 encoder layers, 2 heads,
    $d_{\text{model}} = 32$, feedforward 64, CLS pooling. Parameter
    count matched to the LSTM at $\sim$30k.
    \item \texttt{transformer (big)}: 4 encoder layers, 4 heads,
    $d_{\text{model}} = 128$, feedforward 512, GELU activations,
    pre-norm, sinusoidal positional encoding. Trained with AdamW,
    learning rate $5 \times 10^{-4}$, weight decay $10^{-4}$, cosine
    learning rate with one warmup epoch, gradient clipping at 1.0, and
    validation-based early stopping (patience 5, validation slice
    being the last $10\%$ of each training window). $\sim$3M parameters.
\end{itemize}

\subsection{Evaluation}

We evaluate via rolling walk-forward with 90-day training windows,
30-day test windows, and 30-day step. With $t_0 = 2024\text{-}01\text{-}01$,
180 days of training history before $t_0$, and a 750-day test horizon
(capped at the snapshot date), this yields 27 walk-forward splits and
135{,}000 test samples in total. Each model is retrained from scratch
on each split; deep models use the same seed across splits to limit a
secondary source of variance.

Per-split metrics include accuracy, macro F1, total spread-signal PnL,
the ratio $\text{PnL}_{\text{mean/std}}$, win rate per acting bar, and
the maximum drawdown of the cumulative spread-signal PnL. We compute
block-bootstrap 5\%--95\% confidence intervals over the 27 per-split
values, with block size 3 to reflect the $\sim$3-step training overlap
implied by the 90-day-train, 30-day-step schedule.

\paragraph{HAC correction.} The ratio
$\text{PnL}_{\text{mean/std}}$ uses the bar-level sample standard
deviation as its denominator. Because the target is a 24-hour
forward-looking quantity computed on hourly bars, consecutive samples
share 23 hours of target window and are not independent. We replace the
sample standard deviation with the Newey--West long-run standard
deviation \citep{newey1987hac} at lag 24,
\begin{equation}
    \widehat{\sigma}^2_{\text{NW}} = \widehat{\gamma}_0
        + 2 \sum_{j=1}^{24} \left(1 - \tfrac{j}{25}\right) \widehat{\gamma}_j,
\end{equation}
where $\widehat{\gamma}_j$ is the sample autocovariance at lag $j$.
This correction shrinks the headline ratio by roughly a factor of
$1.7$--$2.0$ across our trading models.

\subsection{Cost-aware backtester}

We implement a portfolio backtester operating on the walk-forward
predictions. The position model is dollar-hedged: a long position
holds $L$ dollars of $A$ and short $L \beta$ dollars of $B$, where
$\beta$ is the static OLS hedge from pair selection (used only for
contract sizing, not for spread computation). Per-bar PnL is
\begin{equation}
    \pi_t = h_{t-1} \cdot L \cdot (r^A_t - \beta r^B_t),
\end{equation}
where $h_{t-1} \in \{-1, 0, +1\}$ is the position held into bar $t$,
set at the close of bar $t-1$ from the classifier's prediction at
$t-1$. This one-bar execution lag is enforced and tested.

The signal-to-position map operates in one of two modes. The
bar-by-bar mode maps each classifier prediction directly to a target
position; this generates many flickering ``trades'' driven by
classifier noise. The state-machine mode is the one we use throughout:
open when flat and the classifier predicts class 0 and
$|\text{spread}_z| \geq z_{\text{in}}$, close when
$|\text{spread}_z| \leq z_{\text{out}}$ or the classifier predicts
class 2 or the spread sign flips. We use $z_{\text{in}} = 2$,
$z_{\text{out}} = 0.5$.

Transaction cost is charged at $(\text{taker fee} + \text{slippage})$
basis points per side, applied to $|h_t - h_{t-1}| \cdot L \cdot (1 + |\beta|)$
in dollars. The annualized Sharpe is computed on bar returns
$\pi_t / C$ where $C = N_{\text{pairs}} \cdot L \cdot 2$ is a fixed
capital base (one earlier version of this calculation used a
time-varying deployed-capital denominator that could produce positive
Sharpe on negative total return; this was corrected). The backtester
passes six hand-computable sanity tests, including a Sharpe-annualization
check, a one-bar-lag injection test, and a state-machine
open-then-close test.

\subsection{Inference for selection bias}

Because we evaluate many strategy configurations (two spread
definitions, two label schemes, seven models, multiple HMM variants,
multiple entry/exit thresholds, three cost levels), the best
observed Sharpe is biased upward by selection. We apply the deflated
Sharpe correction of \citet{bailey2014dsr}. Given the variance
$\widehat{V}$ of Sharpe estimates across $N$ trial configurations and
the observed Sharpe $\widehat{SR}$ of a candidate strategy with
$T$ return observations, skewness $\widehat{\gamma}_3$ and kurtosis
$\widehat{\gamma}_4$,
\begin{equation}
    \widehat{SR}_0 = \sqrt{\widehat{V}} \cdot \left[
        (1 - \gamma_e) \Phi^{-1}\left(1 - \tfrac{1}{N}\right)
        + \gamma_e \Phi^{-1}\left(1 - \tfrac{1}{N e}\right) \right],
\end{equation}
\begin{equation}
    \text{DSR}(\widehat{SR}) = \Phi\left(
        \frac{(\widehat{SR} - \widehat{SR}_0)\sqrt{T-1}}
             {\sqrt{1 - \widehat{\gamma}_3 \widehat{SR}
                     + \tfrac{\widehat{\gamma}_4 - 1}{4} \widehat{SR}^2}}\right),
\end{equation}
where $\gamma_e \approx 0.5772$ is the Euler--Mascheroni constant.
$\text{DSR}$ is the probability that the candidate's true Sharpe
exceeds zero after correction. We report DSR for $N \in \{5, 10, 25,
50\}$, with $N = 50$ as the conservative estimate of trial count
across this project.

\subsection{Audit and leakage tests} \label{sec:audit}

The pipeline is verified by 10 leakage tests
(\texttt{tests/test\_leakage\_audit.py}) and 6 backtester sanity tests
(\texttt{tests/test\_backtest\_sanity.py}). The leakage tests inject
known future signals at specified times and verify that pre-injection
features, targets, labels, and backtester positions are unchanged.

Two real bugs were caught and fixed during the audit:

\begin{enumerate}
    \item The training-period Kalman residuals were computed with
    default filter parameters while the test-period residuals used
    fitted parameters, creating a parameter discontinuity at the
    train/test boundary. The fix returns training residuals from
    \texttt{fit\_kalman\_mle} computed with fitted parameters during
    the same forward pass that records the final state.

    \item The per-pair label threshold was fit on the first $50\%$ of
    each pair's data when the per-pair-label option was set. With our
    long pipeline (180-day train, 750-day test) the first $50\%$
    extends $\sim$9 months past $t_0$, leaking future spread vol into
    the label distribution. The fix passes an explicit
    \texttt{label\_train\_end} timestamp ($= t_0$) to the dataset
    builder.
\end{enumerate}

Headline numbers in the following section are computed on the
post-audit pipeline. The Kalman cointegration finding
(Section~\ref{sec:finding1}) is unaffected by the audit fixes because
its OOS column never used the buggy code path.

\section{Results} \label{sec:results}

\subsection{Static OLS fails out-of-sample; Kalman recovers cointegration} \label{sec:finding1}

For each of the top 10 correlation-ranked pairs in the
liquidity-filtered universe, we fit (i) the static OLS hedge ratio on
a 90-day training slice, and (ii) the Kalman filter
(Section~\ref{sec:kalman}) by MLE on the same training slice, then
forward-roll on a 30-day held-out test slice with parameters and
trained final state held fixed. We ADF-test both residual series on
the test slice.

\begin{table}[h]
\centering
\caption{Out-of-sample ADF $p$-values, static OLS versus Kalman, 10 pairs.}
\label{tab:oos-10}
\small
\begin{tabular}{lrrr}
\toprule
Pair & static OOS $p$ & Kalman OOS $p$ & Kalman $q_\beta$ \\
\midrule
DOGEUSDT\_SUIUSDT  & 0.923 & $1.5 \times 10^{-7}$ & $7.4 \times 10^{-6}$ \\
SOLUSDT\_DOGEUSDT  & 0.985 & $2.2 \times 10^{-7}$ & $4.6 \times 10^{-6}$ \\
XRPUSDT\_SOLUSDT   & 0.373 & $3.0 \times 10^{-7}$ & $6.8 \times 10^{-8}$ \\
XRPUSDT\_DOGEUSDT  & 0.092 & $2.5 \times 10^{-6}$ & $3.5 \times 10^{-6}$ \\
ADAUSDT\_HBARUSDT  & 0.094 & $7.3 \times 10^{-6}$ & $3.9 \times 10^{-6}$ \\
SOLUSDT\_ADAUSDT   & 0.823 & $2.9 \times 10^{-5}$ & $2.5 \times 10^{-5}$ \\
ETHUSDT\_ADAUSDT   & 0.411 & $3.4 \times 10^{-5}$ & $2.9 \times 10^{-5}$ \\
XRPUSDT\_ADAUSDT   & 0.011 & $1.5 \times 10^{-4}$ & $2.6 \times 10^{-5}$ \\
XRPUSDT\_ETHUSDT   & 0.029 & $2.1 \times 10^{-4}$ & $2.8 \times 10^{-5}$ \\
ETHUSDT\_SOLUSDT   & 0.013 & $8.0 \times 10^{-4}$ & $2.0 \times 10^{-5}$ \\
\bottomrule
\end{tabular}
\end{table}

Static OLS holds cointegration at $p < 0.05$ out-of-sample on $2$ of
$10$ pairs. Kalman holds it on $10$ of $10$ at $p < 10^{-3}$ on the
worst case. The MLE-fitted $q_\beta$ lands in $[10^{-8}, 3 \times 10^{-5}]$,
indicating a very slow random walk on $\beta$. This slow drift is the
qualitative check that the filter is not trivially whitening: a fast
random walk on $\beta$ could fit residuals to anything, but a $q_\beta$
of order $10^{-5}$ implies $\beta$ moves by less than $0.005$ per
hour in standard deviation.

\paragraph{Generalization to the full universe.} The same protocol on
all $\binom{50}{2} = 1{,}225$ pairs in the liquidity universe (1{,}219
after dropping pairs with insufficient training overlap) is summarized
in Table~\ref{tab:oos-screen}. The distribution of OOS $p$-values is
shown in Figure~\ref{fig:screen-cdf}.

\begin{table}[h]
\centering
\caption{Universe-wide OOS cointegration counts ($n = 1{,}219$ pairs).}
\label{tab:oos-screen}
\small
\begin{tabular}{lrrrr}
\toprule
$p$ threshold & static count & static \% & Kalman count & Kalman \% \\
\midrule
$p < 0.001$ & 0    & $0.0\%$  & 1{,}174 & $96.3\%$ \\
$p < 0.005$ & 1    & $0.1\%$  & 1{,}188 & $97.5\%$ \\
$p < 0.010$ & 10   & $0.8\%$  & 1{,}198 & $98.3\%$ \\
$p < 0.050$ & 66   & $5.4\%$  & 1{,}215 & $99.7\%$ \\
$p < 0.100$ & 142  & $11.6\%$ & 1{,}217 & $99.8\%$ \\
\bottomrule
\end{tabular}
\end{table}

\begin{figure}[h]
\centering
\includegraphics[width=0.85\linewidth]{figures/kalman_screen_cdf.png}
\caption{CDF of out-of-sample ADF $p$-values, static OLS versus Kalman
dynamic hedge, across $n = 1{,}219$ pairs in the top-50 liquidity
universe at $t_0 = 2024\text{-}01\text{-}01$. Log $p$-axis. Static
OLS clears $p < 0.05$ on only $5.4\%$ of pairs; Kalman clears it on
$99.7\%$.}
\label{fig:screen-cdf}
\end{figure}

This is the load-bearing finding of the paper. The collective evidence
from 1{,}219 independent pair tests is overwhelming: a static OLS hedge
ratio fit on $90$ days of training data does not characterize the
relationship out-of-sample, while a Kalman filter with the same
training does. The cryptocurrency market admits cointegrated pair
relationships, but the relationships are non-stationary and any
analysis that treats $\beta$ as fixed will conclude there is no
cointegration when there is.

\subsection{Walk-forward classifier evaluation}

We evaluate seven models across 27 walk-forward splits on the
post-audit Kalman pipeline with per-pair labels. Block-bootstrap
$5\%$--$95\%$ confidence intervals on $\text{PnL}_{\text{mean/std}}$
and on per-trade win rate are reported in Table~\ref{tab:wf-deep}.

\begin{table}[h]
\centering
\caption{Walk-forward classifier metrics, 27 splits, block-bootstrap CIs.}
\label{tab:wf-deep}
\small
\begin{tabular}{lrrrrr}
\toprule
Model & $\text{PnL}_{\text{mean/std}}$ & 90\% CI & win rate & win rate CI & trades/split \\
\midrule
\textbf{lstm}                & \textbf{0.376} & $[0.362, 0.388]$ & 0.795 & $[0.768, 0.824]$ & 2{,}083 \\
booster                       & 0.356 & $[0.342, 0.368]$ & 0.741 & $[0.720, 0.760]$ & 2{,}355 \\
zscore\_rule                  & 0.282 & $[0.274, 0.291]$ & \textbf{0.963} & $[0.959, 0.966]$ & 588 \\
majority\_class               & 0.021 & $[0.000, 0.063]$ & 0.028 & $[0.000, 0.083]$ & 185 \\
transformer (small)           & 0.018 & $[-0.025, 0.064]$ & 0.437 & $[0.354, 0.519]$ & 691 \\
random\_stratified            & 0.001 & $[-0.003, 0.005]$ & 0.502 & $[0.500, 0.505]$ & 2{,}689 \\
persist\_class                & 0.000 & $[0.000, 0.000]$ & 0.000 & $[0.000, 0.000]$ & 0 \\
\bottomrule
\end{tabular}
\end{table}

LSTM leads the ranking on $\text{PnL}_{\text{mean/std}}$, narrowly
above booster; the two CIs overlap slightly but LSTM's per-trade win
rate is strictly above booster's. The classical $z$-threshold rule is
materially below the ML models on $\text{PnL}_{\text{mean/std}}$ but
strictly above them on per-trade win rate, reflecting a different
operating point: $z$-score acts only on the $\sim$25\% of bars where
$|z| \geq 1.5$ and is correct on $96\%$ of those decisions; the ML
models act on a larger fraction of bars and are correct $74\%$--$80\%$
of those.

The small transformer (parameter-matched to LSTM) is statistically
indistinguishable from random in this evaluation. A larger
transformer ($\sim$3M parameters, cosine LR, validation-based early
stopping), evaluated on a related dataset (per-pair labels without
the audit fixes), reaches $\text{PnL}_{\text{mean/std}} = 0.334$ and
win rate $0.845$ across 27 walk-forward splits; we therefore attribute
the small transformer's failure to insufficient model capacity and
training budget rather than to a property of the transformer
architecture.

\paragraph{HAC correction.} Table~\ref{tab:hac} reports the Newey--West
HAC-corrected $\text{PnL}_{\text{mean/std}}$ alongside the iid
sample-std version. The HAC correction shrinks every trading model by
approximately a factor of $1.7$--$2.1$; the random-prediction inflation
factor of $1.00$ is the sanity check that random predictions do not
exhibit lag autocorrelation.

\begin{table}[h]
\centering
\caption{HAC-corrected $\text{PnL}_{\text{mean/std}}$, Newey--West lag $24$.}
\label{tab:hac}
\small
\begin{tabular}{lrrr}
\toprule
Model & iid & HAC (lag 24) & inflation \\
\midrule
lstm                          & 0.376 & 0.197 & 1.92 \\
booster                       & 0.356 & 0.188 & 1.91 \\
zscore\_rule                  & 0.282 & 0.138 & 2.07 \\
majority\_class               & 0.021 & 0.012 & 1.70 \\
transformer (small)           & 0.018 & 0.012 & 1.29 \\
random\_stratified            & 0.001 & 0.001 & 1.00 \\
\bottomrule
\end{tabular}
\end{table}

\subsection{Feature ablation: the signal reduces to one feature}

We perform a one-feature-out ablation: for each of the 14 tabular
features, retrain the booster with that feature removed and report the
change in held-out accuracy, F1, and total spread-signal PnL relative
to the all-features baseline (Table~\ref{tab:ablation}).

\begin{table}[h]
\centering
\caption{One-feature-out ablation on the booster, post-audit pipeline.}
\label{tab:ablation}
\small
\begin{tabular}{lrrr}
\toprule
Feature dropped & $\Delta$ accuracy & $\Delta$ macro F1 & $\Delta$ total PnL \\
\midrule
(baseline, all 14)                                  & 0.587 & 0.575 & 26.22 \\
\textbf{latest\_spread\_z}                          & \textbf{-0.144} & \textbf{-0.190} & \textbf{-20.72} \\
mean\_volume\_ratio                                  & -0.001 & -0.003 & -0.43 \\
all 12 other features (worst case)                  & $\leq \pm 0.005$ & $\leq \pm 0.005$ & $\leq \pm 0.5$ \\
\bottomrule
\end{tabular}
\end{table}

Dropping \texttt{latest\_spread\_z} costs $14$ percentage points of
accuracy, $19$ of macro F1, and $79\%$ of total PnL. Dropping any of
the other $13$ features changes the booster's PnL by less than half a
unit. The classifier reduces, in practice, to a fancy non-linear
threshold on the spread $z$-score, which is precisely the quantity the
classical $z$-rule already uses.

Reading Table~\ref{tab:wf-deep} in light of this ablation: the LSTM
and booster do extract slightly more signal from
$\text{spread}_z$ than the hard $|z| \geq 1.5$ rule does, but the
``slightly'' is on the order of $\text{PnL}_{\text{mean/std}}$ of
$0.07$--$0.09$ (HAC: $0.05$--$0.06$) above the $z$-rule baseline.

\subsection{Cost sensitivity and the deflated-Sharpe verdict}

We feed the walk-forward predictions through the cost-aware backtester
at three transaction-cost levels: $0$ basis points (signal sanity
check), $5$ bps round-trip per leg (approximate Binance.US maker fee
including modest slippage), and $15$ bps round-trip per leg
(approximate Binance.US taker fee plus slippage). Annualized Sharpe
ratios for the top three models are reported in
Table~\ref{tab:cost-naive}.

\begin{table}[h]
\centering
\caption{Naive (uncorrected) annualized Sharpe by transaction cost.}
\label{tab:cost-naive}
\small
\begin{tabular}{lrrr}
\toprule
Round-trip cost & lstm Sharpe & zscore Sharpe & booster Sharpe \\
\midrule
0 bps           & 8.08  & 8.03  & 8.27 \\
5 bps (maker)   & 0.80  & 0.83  & 0.90 \\
15 bps (taker)  & $-14.04$ & $-13.88$ & $-14.17$ \\
\bottomrule
\end{tabular}
\end{table}

Read naively, these numbers tell a clean story: pre-cost the strategy
clears Sharpe $\sim$8; at maker fees it sits just above zero; at taker
fees it loses. This was our original headline. It does not survive
the deflated-Sharpe correction.

Table~\ref{tab:dsr} reports the deflated-Sharpe verdict for the top
model at each cost level for $N \in \{5, 10, 25, 50\}$ candidate
strategies considered. $N = 50$ is our conservative estimate of the
configurations evaluated across this project; $N = 10$ corresponds to
counting only the seven baseline-plus-classifier rows in
Table~\ref{tab:wf-deep} plus three additional variants. The
expected-max-Sharpe under the null grows with $N$.

\begin{table}[h]
\centering
\caption{Deflated Sharpe (Bailey--L\'opez de Prado) for the top model
at each cost level. $\text{SR}_0$ is the expected maximum annualized
Sharpe under the null over $N$ trials.}
\label{tab:dsr}
\small
\begin{tabular}{rrrrr}
\toprule
$N$ trials & $\text{SR}_0$ (ann.) & DSR @ 0 bps & DSR @ 5 bps & DSR @ 15 bps \\
\midrule
$5$   & $4.08$ & $0.9999$ & $0.42$ & $\sim 0$ \\
$10$  & $5.38$ & $0.99$   & $0.38$ & $\sim 0$ \\
$25$  & $6.83$ & $0.84$   & $0.37$ & $\sim 0$ \\
$50$  & $7.78$ & $0.67$   & $0.34$ & $\sim 0$ \\
\bottomrule
\end{tabular}
\end{table}

Three statements:

\paragraph{The 5 bps break-even claim is retracted.} At any plausible
$N$, the observed Sharpe of $\sim$0.9 at 5 bps round-trip is below the
chance-best of $\sim$1.2 at $N = 10$, with DSR $< 0.5$. We cannot
claim the strategy clears Binance.US maker fees on this data.

\paragraph{The pre-cost Sharpe is robust only at very small $N$.} At
$N = 5$--$10$, the observed $\sim$8 Sharpe is significant. At
$N \geq 25$, marginal. At $N = 50$, the DSR is $0.67$, well below
the conventional $0.95$ threshold for rejecting the null. The
configurations evaluated across this project (two spread definitions,
two label schemes, seven models with multiple variants, HMM filter at
two state counts, two entry/exit threshold pairs, three cost levels,
several preprocessing options) plausibly number in the
$30$--$60$ range. The most honest reading: we cannot confidently
distinguish even the pre-cost signal from selection bias on this
dataset alone.

\paragraph{The taker-fee result is unaffected.} At $15$ bps round-trip
per leg the strategy loses deterministically, with DSR $\approx 0$.
No selection correction is needed to confirm a negative result.

\subsection{Where the Sharpe actually comes from} \label{sec:where}

Table~\ref{tab:where} highlights a single row from the pre-cost
Sharpe data that warrants its own discussion: random predictions in
the state machine produce a Sharpe of $4.78$ annualized, more than
half of what the trained classifiers produce. With no model at all.

\begin{table}[h]
\centering
\caption{Pre-cost Sharpe of all models in the state-machine backtester,
sorted descending.}
\label{tab:where}
\small
\begin{tabular}{lrrr}
\toprule
Model & Annualized Sharpe & Trades & Win rate \\
\midrule
booster              & 8.27 & 12{,}634 & 0.621 \\
lstm                 & 8.08 & 12{,}632 & 0.620 \\
zscore\_rule         & 8.03 & 12{,}612 & 0.621 \\
random\_stratified   & 4.78 &  4{,}048 & 0.582 \\
transformer (small)  & 2.51 &     676 & 0.551 \\
majority\_class      & 1.86 &     520 & 0.632 \\
persist\_class       & 0.00 &       0 & 0.000 \\
\bottomrule
\end{tabular}
\end{table}

The mechanism: the state machine opens a position only when
$|\text{spread}_z| \geq 2$ \emph{and} the classifier predicts the
revert class. The classifier's prediction determines \emph{whether}
the entry condition fires, not whether the entry condition is
predictive. On Kalman-derived stationary spreads
(Section~\ref{sec:finding1}), an entry against the spread sign at
$|z| \geq 2$ has positive expected return regardless of the
classifier's input. A random classifier still fires the entry
condition $\sim$33\% of the time on high-$|z|$ bars and extracts a
proportional share of the alpha.

The defensible reading of Table~\ref{tab:where}: the alpha lives in
the $|z| \geq 2$ entry rule applied to Kalman-derived spreads, not in
the choice of classifier. The booster, LSTM, and $z$-rule cluster at
Sharpe $\sim$8 because all three reliably fire the entry condition
when it is available. The classifier choice contributes at most an
additional $\sim$0.3 Sharpe over the $z$-rule baseline. Random
predictions extract roughly half the alpha because they fire the
entry condition stochastically rather than deterministically.

This observation reinforces the deflated-Sharpe verdict. The relevant
selection is not over classifiers, but over the entry condition
(and over the spread definition, where Kalman versus static is a
real choice). The $|z| \geq 2$ rule on a stationary spread is the
strategy; the classifier is decoration.

\subsection{Negative result: HMM regime filter}

We also tested whether a 2-state Gaussian hidden Markov model on
spread-derived features could isolate productive bars. Per pair, we
fit a 2-state HMM on a vector of spread features within the training
slice, identify the ``mean-reverting'' state as the one with lower
spread-difference standard deviation, decode states for the full
pair history via Viterbi, and suppress non-mean-reverting bars to the
flat class. Table~\ref{tab:hmm} shows the result.

\begin{table}[h]
\centering
\caption{HMM regime filter ablation, 27 walk-forward splits.}
\label{tab:hmm}
\small
\begin{tabular}{lrrr}
\toprule
Model & raw total PnL & HMM-filtered total PnL & change \\
\midrule
lstm                          & 334.6 & 142.6 & $-57\%$ \\
booster                       & 321.5 & 132.8 & $-59\%$ \\
zscore\_rule                  & 195.7 &  71.4 & $-63\%$ \\
transformer (small)           &   0.7 &  $-7.0$ & negative \\
majority\_class               &  21.2 &   9.7 & $-54\%$ \\
\bottomrule
\end{tabular}
\end{table}

The HMM suppresses $40$--$55\%$ of bars and removes $50$--$65\%$ of
PnL on every model tested. The negative result is robust across the
spread definition (static and Kalman), $n_{\text{states}} \in \{2, 3\}$,
and $n_{\text{init}} \in \{1, 3\}$ random restarts with
convergence-preferred selection. A 2-state Gaussian HMM on
$[\text{spread}_z, \Delta\text{spread}, \text{spread vol}]$ does not
recover meaningful mean-reverting regimes on this universe. A
different model class or a different feature combination may; the
configurations tested do not.

\section{Limitations} \label{sec:limits}

\begin{itemize}
    \item \emph{Single exchange.} All findings are specific to
    Binance.US. The cost ceiling implicitly assumes Binance.US fees;
    other US exchanges (Coinbase, Kraken) have different fee
    structures and would shift the break-even.

    \item \emph{Single quote currency.} USDT pairs only. USD-quoted
    pairs would have a different liquidity profile.

    \item \emph{Single bar interval.} Hourly. Lower-frequency alpha
    (4-hour, daily) may have different cost characteristics and is
    not tested here.

    \item \emph{Survivorship.} Quantified at a loose upper bound of
    $6.7\%$ (14 delistings in window over a 209-pair denominator) and
    a realistic upper bound of $\sim$5\% after the liquidity filter.
    The bias direction on the headline metrics is upward.

    \item \emph{No funding cost on the synthetic short leg.} Spot
    USDT shorting on Binance.US in practice requires margin or
    perpetual futures with their own funding rates. Our model treats
    the short leg as costless to hold.

    \item \emph{Flat slippage, no market impact.} Until L2 data
    lands, slippage is modeled as a flat basis-point cost rather
    than as a function of trade size.

    \item \emph{$t_0$ sensitivity not tested.} All results use
    $t_0 = 2024\text{-}01\text{-}01$. Repeating the analysis at other
    $t_0$ values is straightforward and on the next-step list.

    \item \emph{3-class label is a heuristic.} A direct regression
    target on the continuous 24-hour spread change is an obvious
    alternative and is not run.

    \item \emph{Selection bias is real and not fully controlled.}
    The deflated-Sharpe correction depends on the variance of Sharpes
    across the trial set, which we estimate from the 7 models in
    Table~\ref{tab:wf-deep}. A more conservative estimate would
    include all configurations explored during development, which we
    did not log exhaustively.
\end{itemize}

\section{Phase 2: L2 microstructure} \label{sec:phase2}

This work is the first phase of a two-phase project. Phase 2 will
incorporate Level-2 order-book data from Tardis, with three concrete
goals:

\begin{enumerate}
    \item \emph{Microstructure features.} Microprice, quoted-spread in
    basis points, order-book imbalance, and trade-flow signals replace
    the hourly OHLCV summary as the input to the classifier.
    Implementation scaffold is in \texttt{src/l2\_store.py}; the
    parquet schema and feature functions are written and tested on
    synthetic data, with the ingestion parsers left as stubs pending
    the actual Tardis data format.

    \item \emph{Market impact in the backtester.} Replace the flat
    slippage model with a model whose cost depends on instantaneous
    L2 depth and trade size.

    \item \emph{Pre-registered strategy.} Address the selection-bias
    concern (Section~\ref{sec:results}, Table~\ref{tab:dsr}) by
    pre-committing to a single model configuration based on Phase 1
    findings, then evaluating it once on Phase 2 features. This is
    the only fully honest way to defang the deflated-Sharpe critique
    of the headline cost-ceiling claim.
\end{enumerate}

\section{Conclusion}

The Phase 1 evidence supports one strong claim, one weaker claim,
and one negative result.

\paragraph{Strong.} Cryptocurrency pair relationships are real but
non-stationary on multi-month horizons. A static OLS hedge ratio fit
on $90$ days fails to maintain cointegration on $30$-day held-out
test slices for $94.6\%$ of pairs in the top-50 liquidity universe.
A Kalman filter with parameters MLE-fit on the same training and
forward-rolled on test recovers cointegration on $99.7\%$ of the
same pairs. Any analysis of crypto pair-trading that treats $\beta$
as fixed will systematically conclude that no cointegration exists
when it does.

\paragraph{Weaker.} Once a stationary spread is available, a
classical $z$-threshold rule at $|z| \geq 2$ captures most of the
available reversion signal. Modern machine-learning models (gradient
boosting, LSTM, properly-sized transformer) provide a small
additional edge that is statistically detectable on 27 walk-forward
splits but does not survive a Bailey--L\'opez de Prado deflated-Sharpe
correction at realistic transaction-cost levels for the number of
configurations evaluated. The classifier's role is to fire the
$|z| \geq 2$ entry condition; the entry condition itself is the
alpha source.

\paragraph{Negative.} A 2-state Gaussian HMM on spread features does
not isolate productive bars; the filtered version of any model
performs worse than the unfiltered version by $50$--$65\%$. The
negative result is robust to state count, random-start count, and
spread definition.

The market-efficiency statement we can defend: at hourly cadence on
Binance.US USDT pair spreads, the mean-reverting alpha extractable by
the strategies tested is not large enough to clear realistic
transaction costs after honest selection-bias adjustment. Whether
intra-bar execution on L2 data changes this picture is the central
Phase 2 question.

\paragraph{Reproducibility.} All numbers in this paper are produced
by scripts in the project repository at
\url{https://github.com/samsam324/math-199-research}. Command-line
recipes are listed in the repository \texttt{README.md}. The
backtester is verified by 6 hand-computable sanity tests; the
pipeline is verified by 10 leakage tests.

\bibliographystyle{plainnat}
\begin{thebibliography}{99}

\bibitem[Bailey and L\'opez de Prado(2014)]{bailey2014dsr}
Bailey, D.~H. and L\'opez de Prado, M. (2014).
\newblock The deflated Sharpe ratio: Correcting for selection bias,
backtest overfitting, and non-normality.
\newblock \emph{Journal of Portfolio Management}, 40(5):94--107.

\bibitem[Gatev et~al.(2006)Gatev, Goetzmann, and
Rouwenhorst]{gatev2006pairs}
Gatev, E., Goetzmann, W.~N., and Rouwenhorst, K.~G. (2006).
\newblock Pairs trading: Performance of a relative-value arbitrage rule.
\newblock \emph{Review of Financial Studies}, 19(3):797--827.

\bibitem[Newey and West(1987)]{newey1987hac}
Newey, W.~K. and West, K.~D. (1987).
\newblock A simple, positive semi-definite, heteroskedasticity and
autocorrelation consistent covariance matrix.
\newblock \emph{Econometrica}, 55(3):703--708.

\end{thebibliography}

\end{document}
